\section{引言：生成式建模的时代}

我们正处于生成式 AI（Generative AI）的巨大变革时代。生成式建模的核心目标不仅仅是从数据中识别模式并据此做出决策，更重要的是能够\textbf{基于所学习到的模式生成全新的数据实例}。这个看似简单的想法在近年来已经在图像生成、音频生成、自然语言处理等众多领域引发了革命性的变化。

\begin{figure}[H]
\centering
\includegraphics[width=0.85\textwidth]{slides-images/slide-02.jpg}
\caption{哪张脸是真实的？事实上三张都是 AI 生成的假脸\protect\footnotemark}
\end{figure}
\footnotetext{Slides 第2页。来源：thispersondoesnotexist.com}

课程以一个引人深思的例子开始：展示了三张人脸照片，让观众判断哪张是真实的。答案揭示后，三张都是由生成模型合成的假脸。这生动地说明了现代生成模型的强大能力，也引出了本节课的核心议题——理解和掌握深度生成建模的基础理论。

\begin{importantbox}{生成式建模的核心目标}
给定来自某个分布的训练样本，学习一个能够表示该分布的模型 $P_{\text{model}}(x)$，使其尽可能接近真实数据分布 $P_{\text{data}}(x)$。生成模型可用于两大任务：
\begin{itemize}
    \item \textbf{密度估计}（Density Estimation）：学习数据的底层概率分布
    \item \textbf{样本生成}（Sample Generation）：从学到的分布中采样以生成新数据
\end{itemize}
\end{importantbox}

\subsection{监督学习 vs. 无监督学习}

在课程之前的内容中，我们主要关注\textbf{监督学习}（Supervised Learning）问题：给定带标签的数据 $(x, y)$，目标是学习从输入 $x$ 到标签 $y$ 的映射函数。而生成式建模属于\textbf{无监督学习}（Unsupervised Learning）的范畴：我们只有数据 $x$，没有标签，目标是发现数据中隐藏的底层结构。

\begin{figure}[H]
\centering
\includegraphics[width=0.85\textwidth]{slides-images/slide-04.jpg}
\caption{监督学习与无监督学习的对比\protect\footnotemark}
\end{figure}
\footnotetext{Slides 第4页。}

\begin{table}[H]
\centering
\caption{监督学习与无监督学习的对比}
\begin{tabular}{lll}
\toprule
\textbf{维度} & \textbf{监督学习} & \textbf{无监督学习} \\
\midrule
数据 & $(x, y)$，$x$ 为数据，$y$ 为标签 & 仅有 $x$，无标签 \\
目标 & 学习 $x \to y$ 的映射函数 & 发现数据的隐藏/底层结构 \\
典型任务 & 分类、回归、目标检测 & 聚类、特征提取、降维 \\
\bottomrule
\end{tabular}
\end{table}

\subsection{生成式建模的核心原理}

\begin{figure}[H]
\centering
\includegraphics[width=0.85\textwidth]{slides-images/slide-05.jpg}
\caption{生成式建模：从密度估计到样本生成\protect\footnotemark}
\end{figure}
\footnotetext{Slides 第5页。}

生成式建模的根本问题可以表述为：如何学习一个概率分布 $P_{\text{model}}(x)$，使其尽可能接近真实数据分布 $P_{\text{data}}(x)$？

$$P_{\text{model}}(x) \approx P_{\text{data}}(x)$$

一旦我们成功学到了这样的概率分布模型，就可以通过从中采样来生成全新的、从未见过的数据实例。生成的样本将反映训练数据的分布特征。

\subsection{为什么关心生成模型？}

生成模型在现实世界中有三大核心应用场景：

\textbf{1. 去偏差（Debiasing）}：生成模型能够自动发现数据集中的底层特征分布。例如，面对一个人脸数据集，模型可以发现肤色、姿态、光照等特征是否存在过度或不足的代表，从而帮助创建更公平、更具代表性的数据集。

\begin{figure}[H]
\centering
\includegraphics[width=0.85\textwidth]{slides-images/slide-06.jpg}
\caption{利用生成模型发现数据集偏差\protect\footnotemark}
\end{figure}
\footnotetext{Slides 第6页。来源：Amini/Soleymani+ AAAI/AIES 2019。}

\textbf{2. 异常检测（Outlier Detection）}：在自动驾驶等安全关键场景中，生成模型可以通过学习数据的概率分布来识别位于分布尾部的罕见事件（如极端天气、行人），从而检测这些异常情况并调整系统行为。

\begin{figure}[H]
\centering
\includegraphics[width=0.85\textwidth]{slides-images/slide-07.jpg}
\caption{生成模型用于自动驾驶中的异常检测\protect\footnotemark}
\end{figure}
\footnotetext{Slides 第7页。来源：Amini+ IROS 2018。}

\textbf{3. 样本生成（Sample Generation）}：这是最广为人知的应用——利用生成模型创建全新的数据实例。这是 ChatGPT、图像生成、视频生成等 Generative AI 应用的基础。

\begin{figure}[H]
\centering
\includegraphics[width=0.85\textwidth]{slides-images/slide-08.jpg}
\caption{生成模型驱动的各类应用：自然语言、图像视频、生物学\protect\footnotemark}
\end{figure}
\footnotetext{Slides 第8页。}

\subsection{本章小结}

本节引入了生成式建模的基本概念，阐明了从监督学习到无监督学习的转变。生成模型的核心目标是学习数据的概率分布，并可用于密度估计、异常检测和样本生成等场景。本讲将重点介绍两大基础架构：\textbf{Variational Autoencoder (VAE)} 和 \textbf{Generative Adversarial Network (GAN)}，它们都属于 \textbf{Latent Variable Model}（隐变量模型）。


